\subsection{Order of contact of two functions at a point}

1.1.1. Let X be a topological space and \(\theta\) a positive numerical function defined in a neighborhood of a point \(x_0\) of X. A function \(f\) , defined in a neighborhood of \(x_0\) and taking values in F, is said to be negligible with respect to \(\theta\) at \(x_0\) if the following condition is satisfied:

For every \(\varepsilon > 0\) and for every continuous seminorm \(\gamma\) on \(F\) , there exists a neighborhood \(V\) of \(x_0\) on which \(f\) and \(\theta\) are defined and such that

\[
\| f (x) \| _ {\gamma} \leqslant \varepsilon \theta (x) \quad \text { for   every } x \in \mathbf {V}.
\]

For \(f\) to be negligible with respect to \(\theta\) , it suffices that this condition be satisfied for a family of seminorms \(\gamma\) defining the topology of F. The fact that \(f\) is negligible or not with respect to \(\theta\) at \(x_{0}\) depends only on the germs of \(f\) and of \(\theta\) at \(x_{0}\) . One denotes by \(o_{x_{0}}(\theta)\) (or \(o(\theta)\) when there is no ambiguity about \(x_{0}\) ) the set of germs at \(x_{0}\) of functions negligible with respect to \(\theta\) at \(x_{0}\) : it is a vector subspace of the space of germs at \(x_{0}\) of mappings with values in F. If \(f\) is negligible with respect to \(\theta\) , we shall write, by abuse of notation, \(f \in o_{x_{0}}(\theta)\) or also \(f(x) \in o(\theta(x))\) when \(x\) tends to \(x_{0}\) .

If \(f\) and \(g\) are two mappings from a neighborhood of \(x_{0}\) to \(F\) , we shall also write \(f \equiv g \mod o(\theta)\) if \(f - g\) is negligible with respect to \(\theta\) .

Suppose that K is equal to R or C and that \(x_{0}\) is adherent to the set Y of points of X where \(\theta\) is defined and nonzero. The relation \(f \in o(\theta)\) then means that \(f(x)/\theta(x)\) tends to 0 as x tends to \(x_{0}\) while remaining in Y, and that \(\theta(x) = 0\) implies \(f(x) = 0\) .

1.1.2. Let \(f\) and \(g\) be two functions with values in \(F\) , defined in a neighborhood of a point \(x_0\) of \(E\) . If \(m\) is a positive integer, we say that \(f\) and \(g\) have contact of order \(\geqslant m\) at \(x_0\) if we have:

\[
f (x) - g (x) \in o (\| x - x _ {0} \| ^ {m}) \quad \text { for } x \text { tending   toward } x _ {0}
\]

whatever norm is chosen to define the topology of E. For this, it suffices that the preceding relation be verified for one norm defining the topology of E. If this is so, we have \(f(x_0) = g(x_0)\) .

If \(f\) and \(g\) take the same value at \(x_{0}\) , we call the order of contact of \(f\) and \(g\) at \(x_{0}\) the upper bound (finite or equal to \(+\infty\) ) of the integers \(m\) such that \(f\) and \(g\) have contact of order \(\geqslant m\) at \(x_{0}\) .

1.1.3. The order of contact of \(f\) and \(g\) at \(x_{0}\) depends only on the germs of these functions at the point \(x_{0}\) . One can therefore speak of the order of contact of two germs \(\varphi\) and \(\psi\) of maps from E to F at the point \(x_{0}\) . The relation \(\langle\varphi\) and \(\psi\) have contact of order \(\geqslant m\rangle\) is an equivalence relation compatible with the vector structure.

